\section{Exact quotient action on the compiled census}
\label{sec:quotients}

\subsection{Full-port coning and its incidence graph}

For $A\subseteq\{0,\ldots,m-1\}$, write $J_A=\bigcup_{j\in A}J_j$.
The operation $\Cone(A)$ replaces $E$ by the equivalence closure of
$E\cup(J_A\times J_A)$. Empty $A$ has no effect. A singleton atom set $A$
need not be a no-op: it identifies all points within that one atom.

Suppose a collection $\mathcal A$ of cone groups is applied. Construct a
graph with one vertex for each profile type and one vertex for each active
atom $j\in\bigcup\mathcal A$. Join a type $\profile$ to an active atom $j$
when $a_j>0$. Join the atom vertices within each $A\in\mathcal A$, using a
star or path. Let $\mathcal K$ denote the connected components of this graph
that contain a type vertex, and let $T$ be the set of its type vertices.

\begin{theorem}[Multiplicity-aware quotient formula]
\label{thm:quotient}
The profile census after applying $\mathcal A$ is obtained as follows:
\begin{enumerate}[label=(\roman*)]
\item Every untouched type $\profile\notin T$ retains its original
multiplicity $h_E(\profile)$.
\item Each graph component $K\in\mathcal K$ contributes exactly one orbit,
with profile
\[
 b(K)=\sum_{\profile\in K}h_E(\profile)\profile.
\]
\item Equal resulting profiles are coalesced by adding multiplicities.
\end{enumerate}
In particular the orbit count is
\begin{equation}
 C_{\mathcal A}=\sum_{\profile\notin T}h_E(\profile)+|\mathcal K|.
 \label{eq:quotient-count}
\end{equation}
The resulting census has at most $s$ distinct rows.
\end{theorem}

\begin{proof}
First consider a touched type $\profile$ and an incident active atom $j$.
Each of its $h_E(\profile)$ original orbits contains a point in $J_j$.
Some applied cone contains the entire atom, so all those points, and
therefore all copies of that type, become equivalent. This justifies using
one type vertex instead of $h_E(\profile)$ vertices. It also explains why
the factor $h_E(\profile)$ must occur in $b(K)$.

An edge from a type to an atom means that the corresponding orbit meets the
coned set. Edges within a cone group express the same full identification.
Thus any two type vertices in one graph component become a single orbit.
Conversely, an identification path between original orbits alternates old
orbit connections and cone connections. Replacing its old orbit occurrences
by their types and its cone occurrences by their atom vertices gives a path
in this graph. No two distinct graph components can therefore merge. An
untouched original orbit meets no coned set and retains its identity.

All original points in the copies of a touched type enter the same new
orbit. Their atom counts add, giving $b(K)$. This proves the census formula
and \eqref{eq:quotient-count}. If $t$ types are touched, at most $t$ graph
components contain types, while at most $s-t$ untouched rows remain. Merging
equal outputs can only reduce this number, proving the row bound.
\end{proof}

No extra orbit is assigned to a purely virtual atom component. Because every
atom is nonempty and the profile table covers its mass, every active atom
in a valid complete table actually meets a type. The explicit condition in
the theorem remains useful in the implementation and its checker.

\begin{example}[A binary population collapses once]
Let every atom have length $q$, and suppose the source consists of $q$
parallel components, one point in each atom. Then
$h_E((1,\ldots,1))=q$. Coning any single atom produces one component with
profile $(q,\ldots,q)$, not $q$ components and not one component of profile
$(1,\ldots,1)$. The transformation uses one multiplication per coordinate;
there is no loop of length $q$.
\end{example}

\subsection{The exact information retained by the table}

\begin{theorem}[Relative completeness and minimality]
\label{thm:observational}
For two equivalence relations $E,F$ on the same partitioned set $X$, the
following conditions are equivalent:
\begin{enumerate}[label=(\roman*)]
\item $h_E=h_F$.
\item There exists a bijection of $X$ preserving each atom setwise and
sending the $E$-orbits bijectively to the $F$-orbits.
\item Every finite sequence of full-port cones, followed by any
atom-constant additive observation, has the same component-weight histogram
for the two sources.
\end{enumerate}
Thus the census is a complete invariant for this observer interface, and
any exact summary supporting every such observation must distinguish
unequal censuses.
\end{theorem}

\begin{proof}
If the censuses agree, pair the orbits of each profile type. For each paired
orbit and each atom, choose a bijection between their equally sized
intersections. Their union is a bijection of $X$ with the required
properties. Conversely such a bijection preserves every orbit profile,
proving (i)$\Leftrightarrow$(ii).

Every full-port cone is unchanged by an atom-preserving permutation, and an
atom-constant weight is unchanged as well. Therefore the bijection remains
an isomorphism after any cone sequence and preserves its weight histogram.
This proves (ii)$\Rightarrow$(iii). Finally take no cones and use the scalar
weight of \cref{thm:packing}. Its histogram encodes the entire census, so
(iii) implies (i). The last assertion follows by using this same
observation to distinguish any two unequal tables.
\end{proof}

The bijections in this proof may be arbitrary permutations inside atoms;
they are not asserted to be interval isometries or geometric homeomorphisms.
The theorem identifies exactly the information forgotten by the observer.
It does not imply that the full three-dimensional objects represented by
$E$ and $F$ are equivalent.

\subsection{A small obstruction to more general attachments}

\begin{proposition}[Profiles do not determine partial attachment outcomes]
\label{prop:barrier}
A complete source-profile census is insufficient to compute the effect of
an arbitrary additional interval pairing, even when that pairing has
singleton domain and range.
\end{proposition}

\begin{proof}
Let $X=\{0,1,2,3\}$, with atoms $\{0,1\}$ and $\{2,3\}$. Set
\[
 E=\{\{0,2\},\{1,3\}\},\qquad
 F=\{\{0,3\},\{1,2\}\}.
\]
Both have $h((1,1))=2$. Add the singleton pairing $0\sim2$. It is already
present in $E$, leaving two orbits. In $F$ it connects the two original
orbits, leaving one. The input summaries are identical and the requested
additional pairing is identical, but the answers differ.
\end{proof}

A related error would be to refine an atom after compilation and infer its
subatom profiles from the old table. Splitting the first atom in the example
into $\{0\}$ and $\{1\}$ immediately requires information that the old table
forgot. Refining every point would avoid this particular loss, but can
require $N$ atoms, exponential in $B$. This is a genuine representation
barrier, not a reason to abandon richer compressed attachment descriptions.

\subsection{Incremental evaluation without repeated graph construction}

The implementation maintains disjoint-set nodes for the $s$ original types
and the $m$ atom vertices. Atom nodes are initially inactive and virtual.
A type not yet touched represents $h_E(\profile)$ separate real components,
although only one dormant node is stored for it.

When atom $j$ is activated for the first time, scan its adjacent types. If a
type is still inactive, replace its population by one active node and decrease
the count by $h_E(\profile)-1$. Join each type to the atom node. Joining two
roots that both contain real components decreases the count by one; joining
a virtual root to a real root does not. Once all atoms in the cone group
are activated, join their nodes together. Duplicate activations never scan
adjacency again.

\begin{theorem}[Online amortization]
\label{thm:online}
Let
\[
 e=\sum_{\profile\in\supp h_E}|\{j:a_j>0\}|.
\]
After an $O(sm)$ scan of the dense profile coordinates, a monotone sequence
of cone calls with $u$ total supplied atom occurrences uses
\[
 O\bigl((s+m+e+u)\alpha(s+m)+u\log(m+1)\bigr)
\]
index-level operations, including canonicalization of input groups. Each
type population is collapsed at most once and each incidence edge is
scanned at most once. Count arithmetic uses $B$-bit integers. Materializing
a complete quotient census is a separate $O(sm)$ arithmetic operation and
is not hidden inside the set-union bound.
\end{theorem}

\begin{proof}
The activation rule implements exactly the two stages of
\cref{thm:quotient}: collapse all copies of a touched type, then merge
connected touched types. A type cannot already belong to a disjoint-set
component when its population is first activated, because edges involving
that type are inserted only by activation. The ``real root'' indicator
therefore accounts for every decrease exactly once.

Each atom is activated once, so each incident edge is visited once. There
are at most $s$ first type activations. Group processing adds at most $u$
atom-node union attempts. Union by rank and path compression give the
stated amortized bound, with the sorting term paying for canonical atom
sets. The count stays between zero and the initial $C\le N$, so it uses at
most $B$ bits. Census materialization visits each original type and adds
$h_E(\profile)\profile$ to its current root, then coalesces equal rows.
This last step uses vector arithmetic and sorting rather than only index
operations.
\end{proof}

The prototype stores the original profile table immutably. Branching a
threshold search clones the small mutable set-union state but shares that
table and its adjacency lists. It never substitutes a modified census for
the source certificate. This separation both prevents rebinding mistakes
and makes the exact cost of repeated state trials visible.
